\section{Retrieval}
\label{sec:retrieval}

Retrieval decides which passages the language model sees, and therefore bounds
what any answer can say. It has two paths. A question that names a location in
a document is answered by \emph{exact lookup} on the stored labels; any other
question goes to \emph{hybrid search}, which combines semantic and lexical
similarity. The router tries exact lookup first and falls back to hybrid search
when the question names nothing or nothing carries the named labels. Besides
the passages, it returns a note stating whether the lookup was exact, which
subdivision the question asked for, whether that subdivision was found, and,
if not, what the provision contains instead (\cref{sec:structure}).

\subsection{Hybrid search}

\paragraph{Dense ranking.} The question is embedded with the same model as the
passages, and every chunk in scope---one selected document or the whole
collection---is ranked by cosine distance to it.

\paragraph{Lexical ranking.} The same chunks are ranked by Okapi
BM25~\cite{robertson2009bm25} over lower-cased word tokens.

\paragraph{Fusion.} The two rankings are combined by reciprocal rank
fusion~\cite{cormack2009rrf} (\cref{lst:rrf}). Each chunk $d$ scores
\begin{equation}
  \operatorname{RRF}(d) = \sum_{r \in \{\mathrm{dense},\, \mathrm{BM25}\}}
  \frac{1}{k + \operatorname{rank}_r(d)}, \qquad k = 60,
\end{equation}
where $\operatorname{rank}_r(d)$ is the chunk's position in ranking $r$,
counted from 0.

\begin{lstlisting}[float, language=Python, caption={Reciprocal rank fusion
(\texttt{retrieval/hybrid.py}).}, label={lst:rrf}]
def _rrf(rank_list: list[list[str]], k: int = 60) -> dict:
    scores: dict[str, float] = {}
    for ranking in rank_list:
        for rank, _id in enumerate(ranking):
            scores[_id] = scores.get(_id, 0.0) + 1.0 / (k + rank)
    return scores
\end{lstlisting}

\paragraph{Trade-offs.} The two rankings fail in complementary ways. Dense
embeddings capture paraphrase---``who can complain'' finds a provision on the
right to lodge a complaint---but blur exact tokens: numbers, amounts and
defined terms (``Article~92'', ``EUR~15\,000\,000'') contribute little to a
sentence vector. BM25 matches exact tokens but knows nothing of synonyms.
Fusing by rank rather than by score avoids having to reconcile two
incomparable scales, a bounded cosine distance and an unbounded BM25 score,
and needs no tuned weights. The cost is that magnitudes are discarded: a
passage far ahead in one ranking gets no more credit than one narrowly ahead.
The BM25 index is rebuilt over the chunks in scope for each question. Over the
2\,019 chunks of the AI Act, hybrid retrieval takes about 235\,ms in total,
which is small beside generation; a collection orders of magnitude larger
would need a persistent inverted index instead.

\subsection{Exact lookup for named locations}

\paragraph{Parsing the question.} A pattern finds unit words followed by a
number (``article 66'', ``Art.~66'', ``articles no.~5'', ``page 48'',
``section 3.13''). Plain numbers become integers and anything else upper-cased
text, matching how the labels were stored (\texttt{3.7a} matches
\texttt{3.7A}). Some conventions name a unit by another word, so a unit is
also tried under an alternative reading: in EU acts, ``clause 148'' means
recital 148, and ``section 5'' may mean article 5. A question that matches
nothing is checked for a mistyped unit word: any word of four to ten letters
followed by a number is compared with the unit vocabulary, and a similarity
ratio of at least 0.75 accepts it, so ``artilce 66'' is read as article 66.

\paragraph{Building filters.} Each named unit becomes one or more filters on
the stored metadata. A page becomes a range condition, because a chunk may run
over two pages: first on the printed page numbers, then on the PDF's own page
numbers (\cref{lst:page}). When a question names several units, all of them
are combined first, so that ``chapter 3 page 48'' retrieves the passages on
printed page 48 \emph{within} chapter 3. Only if no chunk carries them all is
each unit tried on its own, the most specific first: page, then recital,
clause, rule, article, section, annex, schedule, chapter and part. Trying the
units separately from the outset returns chapter 3's opening passages for that
question, because ``chapter 3'' alone matches first; the combination is what
lets a page be addressed inside a chapter.

\begin{lstlisting}[float, language=Python, caption={Filters for one named unit; a
page matches every chunk that runs over it (\texttt{retrieval/hybrid.py}).},
label={lst:page}]
def _alternatives(field, value):
    if field != "page":
        return [{field: {"$eq": value}}]
    if not isinstance(value, int):
        return []

    return [
        {"$and": [{"printed_page": {"$lte": value}},
                  {"printed_page_end": {"$gte": value}}]},
        {"$and": [{"page": {"$lte": value}},
                  {"page_end": {"$gte": value}}]},
        {"page": {"$eq": value}},
    ]
\end{lstlisting}

\paragraph{Returning the provision.} The first filter that matches returns
every chunk of the provision, up to 200, in reading order---not the
$k$ most similar chunks. A named provision is the answer to ``what does it
say'' in full, and similarity has no role in deciding which of its parts
matter. The subdivisions inside it are handled in \cref{sec:structure}.

\paragraph{Trade-offs.} Exact lookup is deterministic and explainable: the
same question always returns the same passages, and the interface can say
``exact match'' instead of presenting a similarity score as if it measured
correctness. It is also fast, taking 3--6\,ms against roughly 235\,ms for
hybrid search. Its weakness is that it depends on the question naming units in
a form the parser recognises. A form it does not recognise falls through to
hybrid search, so the worst case is the behaviour of an ordinary RAG system,
never an empty answer.

\subsection{Deciding how many passages}

Early versions let the user choose how many passages to retrieve, between 2
and 12. A fixed number fits no question well: it truncates a long provision
that was named in full, and it pads a focused question with passages that
merely share its vocabulary. The number is now decided by the system, from
what the question asked for:

\begin{itemize}
  \item \textbf{A named subdivision} (``paragraph 2 of Article 7''): exactly
  the chunks that carry the requested labels, which for a paragraph includes
  all of its points and sub-points (\cref{sec:structure}).
  \item \textbf{A named provision} (``Article 93''): the whole provision in
  reading order.
  \item \textbf{Any other question}: the candidates that score close to the
  best one. Of the twelve best candidates after fusion, candidate $i$
  (counted from 0) is kept if
  \begin{equation}
    i < 3 \quad\text{or}\quad s_i \geq \max_j s_j - 0.15,
  \end{equation}
  where $s_i$ is its cosine similarity to the question, and at most eight are
  kept (\cref{lst:close}).
\end{itemize}

On every path the passages are then admitted in order while their total length
stays within 2\,000 words. At about 1.4 tokens per word this is roughly 2\,800
tokens, which leaves room in the model's 4\,096-token context window for the
instructions and the answer.

\begin{lstlisting}[float, language=Python, caption={Keeping the passages close to the
best one (\texttt{retrieval/hybrid.py}).}, label={lst:close}]
def _close_to_best(chunks, distances, metas):
    if not distances:
        return chunks, distances, metas

    best = 1 - min(distances)
    keep = []
    for i, distance in enumerate(distances):
        if i < _FEWEST or 1 - distance >= best - _MARGIN:
            keep.append(i)
    keep = keep[:_MOST]

    return ([chunks[i] for i in keep], [distances[i] for i in keep],
            [metas[i] for i in keep])
\end{lstlisting}

\paragraph{Trade-offs.} The margin is relative rather than absolute because
similarities between a question and the passages of one document cluster
tightly: for typical questions on the AI Act the top twelve candidates fell
between 0.57 and 0.86, so an absolute threshold would keep everything for one
question and nothing for the next. The minimum of three protects against a
lone passage cut mid-sentence, and the maximum of eight against a question so
generic that everything scores alike. The margin and bounds were set by
inspecting score distributions on the development corpora rather than learned,
and they fix only the upper limit; the context budget decides the rest.
